% !TeX root = ../../../main.tex
\section{Notation and Reporting Minimums}

\begin{tabularx}{\textwidth}{>{\bfseries\raggedright\arraybackslash}p{0.20\textwidth} X}
\toprule
Symbol & Meaning\\
\midrule
$N$ & Problem size under an explicitly stated encoding\\
$\epsilon$ & Required additive or relative numerical error\\
$\delta$ & Permitted failure probability\\
$Q_L,Q_{\mathrm{phys}}$ & Logical- and physical-qubit counts\\
$N_{2q},D_{2q}$ & Two-qubit gate count and two-qubit depth after routing\\
$N_T,D_T$ & Non-Clifford $T$ count and dependency depth\\
$N_{\mathrm{shots}}$ & Total circuit executions, including mitigation and controls\\
$p_s$ & Per-run probability of satisfying the success condition\\
$d$ & Quantum error-correcting code distance\\
$R_T$ & Aggregate magic-state production rate\\
$B_C$ & Strongest matched classical baseline\\
$\TTS$ & End-to-end time to reach a target success probability\\
\bottomrule
\end{tabularx}

At minimum, a reported result should bind these symbols to a particular input family, software revision, compiler configuration, machine or machine model, calibration window, classical hardware, and accuracy definition. Without those bindings, resource numbers are not reusable evidence.
